\section{モデリングの目から見た検定4；対応のある群の比較}
\subsection{授業内容}
\subsubsection{科目の中でのこのコマの位置づけ}
ここまでBetween計画についてのモデル化を勧めてきたが，ここからはWithinモデルについて考えることにする。Withinモデルは相関係数の考え方と，階層モデルへの入り口として位置付けることができる。

Withinモデルは対応がある場合であり，変数間に相関関係を想定することになる。相関関係のあるデータ生成分布は，多変量分布であり，正規分布を多次元正規分布に拡張する必要がある。

多変量正規分布が必要とするパラメータはベクトルと行列であることから，Stanにおけるベクトル型，行列型など特殊な型についての理解を進める。

% またそれとは違うアプローチとして，個人のデータが反復ごとに変化していくようなモデル化を考え，それを数式やコードに落としていくことを考える。
% この考え方に慣れていくと，より複雑な違いを組み込むことができる。\citet{Kubo2012}にあるように，個体差，植木鉢の差など段階を踏まえてモデルを調整していくプロセス，固定効果と変動効果など用語に注意しながら，階層モデルへ展開していく。
\subsubsection{コマ主題細目}
\begin{description}

	\item[対応のある群] 対応のあるデータというのは，反復測定あるいは測定間に相関関係が想定されるデータでとして考えることができる。データの間に相関関係がある場合，多次元正規分布から出てくるモデルになることから，これを使ったデータの書き方を習得する。まずは多次元正規分布の考え方とその表現方法を理解する。

	\item[IDを持ったデータ構造] 既に識別変数をもった整然データのモデル化については習得している。コード化するときは群の識別IDではなく個人の識別IDを持ったコードを作成することになる。代入された変数がさらに代入されるという入子構造のプログラミングに慣れる。

	\item[個人差と変化量を想定した書き方] モデルを3群以上の比較にすることを考えると，多次元モデルをさらに広げることになるが，別の表現の仕方としてベースラインとそこからの変化量として表現できることになる。この時，ベースラインの散らばりすなわち個人差は，別の分布から出てきていることになる。この表記方法は今後の階層モデルにつながる観点でありことに言及する。

\end{description}
\subsubsection{キーワード}
\begin{itemize}
	\item Within計画
	\item 分散共分散行列
	\item ベクトル型と行列型
	\item 多次元正規分布
	      % \item 階層線形モデル
\end{itemize}
\subsection{授業情報}
\paragraph{コマの展開方法}
講義/遠隔/演習
\subsubsection{予習・復習課題}
\paragraph{予習}対応のあるt検定やWithinモデルなど，伝統的な方法による分析方法について復習しておく。
\paragraph{復習}パラメータリカバリや身の回りのデータを使って，今回のモデルが具体的にどのように使うことができるかを考えておく。